\documentclass[../../../include/open-logic-section]{subfiles}

\begin{document}

\olfileid[id]{his}{bio}{can}

\olsection{Georg Cantor}

\olphoto{cantor-georg}{Georg Cantor}

Sebuah biografi awal Georg Cantor (\textsc{gay}-org
\textsc{kahn}-tor) mengklaim bahwa ia dilahirkan dan ditemukan di sebuah
kapal yang sedang berlayar menuju Sankt Peterburg, Rusia, serta bahwa orang
tuanya tidak diketahui. Namun, kisah ini tidak benar; Cantor lahir di Sankt
Peterburg pada tahun 1845.

Cantor meraih gelar doktor dalam matematika di Universitas Berlin pada tahun
1867. Ia dikenal karena karyanya dalam teori himpunan dan dianggap sebagai
pendiri teori himpunan sebagai disiplin penelitian tersendiri. Ia adalah
orang pertama yang membuktikan bahwa ada himpunan-himpunan tak hingga dengan
ukuran berbeda. Teori-teorinya, terutama teorinya tentang berbagai ketakhinggaan,
memicu banyak perdebatan di kalangan matematikawan pada masa itu, dan
karyanya menuai kontroversi.

Keyakinan agama Cantor dan karya matematikanya terjalin erat; ia bahkan
mengklaim bahwa teori bilangan transfinit disampaikan kepadanya secara
langsung oleh Tuhan. Pada masa akhir hidupnya, Cantor mengalami gangguan
kesehatan mental. Sejak tahun 1894, dan makin sering menjelang akhir
hayatnya, Cantor menjalani perawatan di rumah sakit. Kritik keras terhadap
karyanya, termasuk perselisihannya dengan matematikawan Leopold Kronecker,
menyebabkan depresi dan hilangnya minat terhadap matematika. Selama episode
depresi, Cantor beralih ke filsafat dan sastra, bahkan menerbitkan teori
bahwa Francis Bacon adalah penulis drama-drama Shakespeare.

Cantor meninggal pada 6 Januari 1918 di sebuah sanatorium di Halle.

\begin{reading}
Untuk biografi lengkap Cantor, lihat \citet{Dauben1990} dan
\citet{Grattan-Guinness1971}. Pandangan radikal Cantor juga diuraikan dalam
program BBC Radio 4 \emph{A Brief History of Mathematics}
\citep{Sautoy2014}. Jika Anda ingin mendengar teori-teori Cantor dalam
bentuk rap, lihat \citet{Rose2012}.
\end{reading}

\end{document}
